\subsection{R-20: a five-row certificate reaches the rectangular padding boundary}

\paragraph{Statement R-20.}
Put $M_0=8{,}000{,}001$. For every integer $m\geq M_0$, let
\[
 h_m=\left\lfloor\frac{m-97}{5}\right\rfloor,
 \qquad J_m=8h_m(m-317).
\]
Then
\begin{equation}\label{R20:dimension}
 \dim\mathcal P_{m,4}\geq2m^2+2+J_m.
\end{equation}
Consequently, for every pair of integers
\begin{equation}\label{R20:rectrange}
 m\geq M_0,\qquad m<n<4m/3,
\end{equation}
the explicit rectangular multiplicity gaps of G-11 hold, with its
labels $\lambda=(nd/4)^4$, degrees $d=E_nk$, and $k\geq K_n$.
Together with G-12, this makes $n<4m/3$ the exact eventual range,
among $n>m$, in which a rectangular multiplicity gap exists.
This is a statement about the specified four-row sector and its
large coordinate degrees.

More generally, whenever $m\geq M_0$, $n>m$, and
\begin{equation}\label{R20:allrange}
 2n^2\leq2m^2+J_m+2,
\end{equation}
R-14/G-10 give an exactly-four-row multiplicity gap in every degree
$d\geq d_n$. If also $n\geq4m/3$, every such gap is nonrectangular
by G-12. In particular, for every fixed $4/3<\rho<3/\sqrt5$, it includes
$n=\lfloor\rho m\rfloor$ for all sufficiently large $m$.
No maximal generic-rank formula is assumed.

\begin{proof}
We prove the dimension bound by an explicit differential certificate.

\paragraph{1. One simultaneous local construction.}
Let $q$ be a primitive $m$th root, to be chosen more carefully below,
and put
\[
 \ell_i=x-q^iy,\quad v_i=z-q^iw,\qquad
 L_c=x-cy,\quad V_c=z-cw.
\]
Start with the cyclic pencil of R-15. At a block starting at zero,
replace the five rows $0,\ldots,4$, in columns $0,\ldots,5$, by
\begin{equation}\label{R20:motif}
 \begin{pmatrix}
 \ell_0&v_0&\varepsilon L_2&\varepsilon L_3&\varepsilon V_3&0\\
 0&\ell_1&v_1&\varepsilon L_4&\varepsilon L_3&\varepsilon V_2\\
 0&0&\ell_2&v_2&\varepsilon V_6&0\\
 0&\varepsilon V_5&0&\ell_3&v_3&\varepsilon V_4\\
 0&\varepsilon L_2&\varepsilon V_7&\varepsilon L_5&\ell_4&v_4
 \end{pmatrix}.
\end{equation}
At starting positions $a=0,5,\ldots,5(h_m-1)$ use the translated
motif: shift all row and column indices by $a$ and replace
$(y,w)$ by $(q^ay,q^aw)$ in its entries. The five-row blocks are
disjoint. The four rows $a+1,\ldots,a+4$ have support only in
columns $a+1,\ldots,a+5$. The four columns
$a+1,\ldots,a+4$ have support only in rows $a,\ldots,a+4$.
Both assertions hold with every block present, since a neighboring
block shares only the boundary column. All remaining entries are
those of the original cycle.

Here is the local cofactor calculation used for both supports.
For a $4\times5$ matrix $B(\varepsilon)$ with diagonal entries
$a_i$, superdiagonal entries $b_i$, and other entries
$\varepsilon c_{ij}$, form $\widetilde B$ by multiplying the
off-bidiagonal entry $(i,j)$ by
\[
 \sigma_{ij}=\begin{cases}
 (-1)^{j-i-1},&j>i+1,\\
 (-1)^{i-j},&j<i.
 \end{cases}
\]
Write $P_j=\operatorname{per}(B\text{ with column }j\text{ deleted})$
and $D_j=\det(\widetilde B\text{ with column }j\text{ deleted})$.
There is one matching at order zero. A matching with one additional
edge $(i,j)$ has permutation sign $(-1)^{j-i-1}$ for $j>i+1$,
or $(-1)^{i-j}$ for $j<i$: the forced remaining bidiagonal edges
form the corresponding consecutive cycle. Thus
\begin{equation}\label{R20:secondorder}
 P_j-D_j=\varepsilon^2 E_j+O(\varepsilon^3),
 \qquad E_0=E_4=0.
\end{equation}
For completeness, the entire second-order expansion is
\begin{align*}
 E_1/2={}&c_{02}c_{10}b_2b_3
 +c_{03}(c_{10}a_2b_3+c_{20}b_1b_3)\\
 &+c_{04}(c_{10}a_2a_3+c_{20}b_1a_3+c_{30}b_1b_2),\\
 E_2/2={}&c_{13}(c_{21}a_0b_3+c_{20}b_0b_3)\\
 &+c_{14}(c_{21}a_0a_3+c_{31}a_0b_2
                   +c_{20}b_0a_3+c_{30}b_0b_2)\\
 &+c_{03}c_{20}a_1b_3
 +c_{04}(c_{20}a_1a_3+c_{30}a_1b_2),\\
 E_3/2={}&c_{24}(c_{32}a_0a_1+c_{31}a_0b_1+c_{30}b_0b_1)\\
 &+c_{14}(c_{31}a_0a_2+c_{30}b_0a_2)
 +c_{04}c_{30}a_1a_2.
\end{align*}
These identities can be verified directly by the $24$ permutations
of each $4\times4$ minor: the displayed $6,9,6$ terms are precisely
the two-additional-edge matchings whose signed determinant term is
negative. This specifies the finite calculation completely.

The determinant row relation gives
\begin{equation}\label{R20:relation}
 \sum_{j=0}^4(-1)^{i+j}\widetilde B_{ij}P_j
 =\varepsilon^2(a_iE_i-b_iE_{i+1})+O(\varepsilon^3).
\end{equation}
Apply it first to rows $1,\ldots,4$ and columns $1,\ldots,5$
of \eqref{R20:motif}, and then to the transpose of its submatrix
in rows $0,\ldots,4$ and columns $1,\ldots,4$.
In the global complementary permanents these identities have a
common remaining factor. For an external deleted row $s$, the
first set factors because the four supported rows must occupy the
four remaining supported columns. For an external deleted column
$c$, the second set factors because the four supported columns
must occupy the four remaining supported rows. The presence of
other blocks does not change either argument.

The coefficient of $\varepsilon^2$, divided by two, has the following
highest components in $(z,w)$ degree. This list follows by substituting
\eqref{R20:motif} into the preceding complete expansion:
\begin{align*}
 R_0={}&-v_1v_2(L_3V_5v_4+V_2V_5\ell_4+V_2L_2v_3),\\
 R_1={}&-v_1v_2v_4V_5V_6,\\
 R_2={}&\ell_3v_1v_4V_5V_6
          -v_2v_3V_4(\ell_1V_7+L_2v_1),\\
 R_3={}&\ell_4v_2V_4(\ell_1V_7+L_2v_1),\\[2pt]
 S_0={}&-\ell_1L_2v_2(\ell_4V_5+L_2v_3),\\
 S_1={}&v_1L_2v_2(\ell_4V_5+L_2v_3)
          -\ell_2L_4v_0v_3V_7,\\
 S_2={}&L_4v_0v_2v_3V_7,\\
 S_3={}&v_3(L_5V_6v_0v_1+L_3V_7v_0v_2
                +\ell_1V_7V_3v_2+L_2V_3v_1v_2).
\end{align*}
Their $(z,w)$ degrees are respectively
\begin{equation}\label{R20:grades}
 (4,5,4,3)\quad\hbox{and}\quad(2,3,4,4).
\end{equation}
Every listed polynomial has total degree five. The complete
coefficient in \eqref{R20:relation} may have additional components
of strictly smaller $(z,w)$ degree; those will be handled below,
without projecting an actual differential image by an unjustified
operation.

At $\varepsilon=0$ the common outside factors are exactly
\begin{equation}\label{R20:outside}
 T_s^{\rm row}=\ell_0\prod_{i=5}^{s-1}v_i
                         \prod_{i=s+1}^{m-1}\ell_i,
 \qquad
 T_c^{\rm col}=\prod_{i=5}^{c-1}\ell_i
                         \prod_{i=c}^{m-1}v_i.
\end{equation}
This follows either by the unique cycle matching or by successively
matching the forced boundary column. If a local polynomial has
$(z,w)$ degree $g$, choose
\begin{equation}\label{R20:alignment}
 s=k+5-g\quad\hbox{for a row family},\qquad
 c=m+g-k\quad\hbox{for a column family}.
\end{equation}
Its highest component times \eqref{R20:outside} then has bidegree
$(m-k,k)$. All deleted indices are external and distinct as required
when $9\leq k\leq m-9$.

\paragraph{2. Fixed rational-function matrices controlling all modes.}
Write $K=q^k$, $X=q^t$. The following construction defines two
$6\times6$ matrices over $\mathbb Q(q,K,X)$, using only the eight
displayed degree-five polynomials. It therefore does not depend on
a numerical extrapolation in $m$.

Let $F$ be one such polynomial of bidegree $(a,g)$, $a+g=5$.
For a row polynomial, express it uniquely as
\begin{equation}\label{R20:rowbasis}
 F=\sum_{j=1}^{a+1} Q^{\rm row}_{a,j}(x,y) C_j(z,w),\qquad
 Q^{\rm row}_{a,j}=
 \prod_{u=j}^{a}(x-Kq^uy)\prod_{u=1}^{j-1}(x-q^uy).
\end{equation}
Define the binary quintic
\begin{equation}\label{R20:rowtransform}
 \mathcal R(F)=\sum_{j=1}^{a+1}X^j C_j(z,q^{-j}w)
       \prod_{u=5-j}^{3}(z-q^uw)
       \prod_{u=-1}^{a-1-j}(z-Kq^uw).
\end{equation}
As everywhere here, a product with lower endpoint greater than its
upper endpoint is one. For a column polynomial use
\begin{equation}\label{R20:colbasis}
 \begin{split}
 F&=\sum_{j=0}^{a}Q^{\rm col}_{a,j}(x,y)C_j(z,w),\\
 Q^{\rm col}_{a,j}&=
 \prod_{u=5-j}^{4}(x-q^uy)
 \prod_{u=0}^{a-1-j}(x-K^{-1}q^{5-a+u}y),
 \end{split}
\end{equation}
and define
\begin{equation}\label{R20:coltransform}
 \mathcal S(F)=\sum_{j=0}^{a}X^{-j}C_j(z,Kq^{j-5}w)
       \prod_{u=-a+j}^{-1}(z-q^uw)
       \prod_{u=0}^{j-1}(z-Kq^{-5+u}w).
\end{equation}
The coefficient matrices of the $Q$ bases are generically invertible.
Their nonzero determinants are also included in the certificate below,
so no problematic specialization is silently inverted.

The row matrix consists of coefficient columns, in the order
$z^5,z^4w,\ldots,w^5$, of
\[
 z\prod_{u=0}^3v_u,\quad w\prod_{u=0}^3v_u,
 \quad\mathcal R(R_0),\ldots,\mathcal R(R_3).
\]
The column matrix consists of the corresponding columns of
\[
 z\prod_{u=-5}^{-2}(z-Kq^uw),\quad
 w\prod_{u=-5}^{-2}(z-Kq^uw),
 \quad\mathcal S(S_0),\ldots,\mathcal S(S_3).
\]
At $(q,K,X)=(6,-1,-1)$ over $\mathbb F_{41}$, these matrices are
\begin{equation}\label{R20:matrices}
 \begin{split}
 M_R&=\begin{pmatrix}
 1&0&36&1&21&34\\28&1&24&28&24&4\\
 34&28&3&28&35&22\\21&34&34&6&27&27\\
 39&21&0&33&24&34\\0&39&26&27&22&6
 \end{pmatrix},\\
 M_C&=\begin{pmatrix}
 1&0&30&28&31&1\\2&1&33&10&24&8\\
 19&2&1&9&13&0\\34&19&35&4&27&5\\
 2&34&13&36&15&4\\0&2&32&3&11&27
 \end{pmatrix}.
 \end{split}
\end{equation}
Their determinants are respectively $16$ and $15$ modulo $41$.
The four row interpolation determinants are $29,1,29,38$;
the four column interpolation determinants are $25,6,32,32$.
All entries are obtained by the products and linear systems
\eqref{R20:rowbasis}--\eqref{R20:coltransform}, which specify a
small reproducible exact certificate.

We record explicit degree bounds in order to turn this certificate
into a uniform statement. Let $d_R$ be the determinant of a row
coefficient matrix in \eqref{R20:rowbasis}, and let $d_C$ be the
determinant of the coefficient matrix of
$K^a Q^{\rm col}_{a,j}$ in \eqref{R20:colbasis}.
For all eight forms, $0\leq a\leq3$. Cramer's rule shows that
the coefficients of
\[
 q^{60}X^5d_R\mathcal R(F),\qquad
 q^{60}X^5d_C\mathcal S(F)
\]
are polynomials in $q,K,X$. Each has degree at most
$150,20,10$ in those three variables separately. Here are direct
bounds verifying the clearing and the stated estimates.
Every coefficient of a local form has $q$-degree at most $20$.
A row interpolation entry has $q$-degree at most $a(a+1)/2$
and $K$-degree at most $a$; its adjugate entry has degrees at most
$18,9$. The substitutions and additional products in
\eqref{R20:rowtransform} have negative $q$ exponent no smaller
than $-10$ and positive contribution at most $7$; after clearing,
the $q$-degree is at most $105$, the $K$-degree at most $12$,
and the $X$-degree at most $9$.
For the column basis, multiplying by $K^a$ gives entries of degrees
at most $4a,a$. Its adjugate has degrees at most $36,9$.
The solution acquires an additional factor $K^a$, and
\eqref{R20:coltransform} has negative $q$ exponent no smaller than
$-31$. After clearing its degrees are at most $116,17,5$.
These are bounded by $150,20,10$ as asserted.

Multiply the first two columns of each matrix by $q^{25}$,
and clear each of its four other columns as just described.
The two base columns now have degrees at most $31,4,0$.
Let $\Delta_R,\Delta_C$ be the two resulting polynomial determinants,
and define
\[
 \mathscr F(q,K,X)=\Delta_R\Delta_C
                  \prod_{\text{four row forms}}d_R
                  \prod_{\text{four column forms}}d_C.
\]
Each interpolation determinant has degrees at most $48,12,0$.
Consequently
\begin{equation}\label{R20:degrees}
 \deg_q\mathscr F\leq2000,\qquad
 \deg_K\mathscr F\leq300,\qquad
 \deg_X\mathscr F\leq80.
\end{equation}
For example the sharper bounds from the preceding estimates are
$1708,272,80$. At the certificate specialization all eight
interpolation determinants have product $20$ modulo $41$;
$K^{a(a+1)}=1$ there, so the column clearing does not change their
listed values. Thus
\[
 \mathscr F(6,-1,-1)
 =6^{580}(-1)^{40}\,20^2\,16\,15=22\pmod{41}.
\]
In particular $\mathscr F$ is a nonzero integer polynomial.

\paragraph{3. Why these matrices are the required Fourier quotients.}
Use the cyclic action $\tau:(y,w)\mapsto(qy,qw)$ of R-16 and
write $A(x,y)=\prod_{u=k}^{m-1}\ell_u$.
All its coefficients $A_r$, $0\leq r\leq m-k$, are nonzero by
the geometric-product formula already proved there. In character
$t$, divide the coefficient with $y$-degree $r$ and $w$-degree
$s$ by $A_r$; here $r+s=t$.

For a row form of $(x,y)$ degree $a$, the outside factor is
$\ell_0\prod_{u=k+a+1}^{m-1}\ell_u$ in those variables.
Multiplication by each $Q^{\rm row}_{a,j}$ in
\eqref{R20:rowbasis} completes exactly $A(x,q^jy)$.
Dividing its coefficient by $A_r$ gives $q^{jr}=X^jq^{-js}$.
The remaining $(z,w)$ factors are
$C_j(z,q^{-j}w)\prod_{u=5-j}^{k+a-1-j}v_u$.
Their common product is
\[
 H_R=\prod_{u=4}^{k-2}v_u,
\]
and the residual quintic is precisely
\eqref{R20:rowtransform}.

For a column form, the outside $(x,y)$ factor is
$\prod_{u=5}^{c-1}\ell_u$, with $c=m+5-a-k$.
The $j$th basis term in \eqref{R20:colbasis} completes
$A(x,q^{5-k-j}y)$. After coefficient division and multiplication
by the nonzero character scalar $X^{k-5}$, its factor is
$X^{-j}q^{(k-5+j)s}$.
Shifting the outside $(z,w)$ roots $c,\ldots,m-1$ by
$k-5+j$ gives indices $-a+j,\ldots,k-6+j$ modulo $m$.
Their common product is
\[
 H_C=\prod_{u=0}^{k-6}v_u,
\]
and the residual quintic is exactly
\eqref{R20:coltransform}. This checks both shifts and both
normalizations, including their dependence on $k,t$.

Suppose $9\leq k,t\leq m-9$ and
$\mathscr F(q,q^k,q^t)\ne0$. The two matrices then show that
base plus row columns span $H_R\operatorname{Sym}^5\mathbb C^2$,
and base plus column columns span
$H_C\operatorname{Sym}^5\mathbb C^2$ before restriction to the
allowed coefficient interval. Their common factor is
$H=\prod_{u=4}^{k-6}v_u$. After dividing by $H$, the first space is
\[
 \left(\prod_{u=k-5}^{k-2}v_u\right)
                 \operatorname{Sym}^5\mathbb C^2,
\]
and the second is
\[
 \left(\prod_{u=0}^{3}v_u\right)
                 \operatorname{Sym}^5\mathbb C^2.
\]
The two displayed degree-four factors are coprime: their indices
are distinct and lie below $m$. Their intersection consists exactly
of their product times $\operatorname{Sym}^1\mathbb C^2$,
which is the two-dimensional base space. The joint span is therefore
all $H\operatorname{Sym}^9\mathbb C^2$, of dimension ten.
This proves joint independence; separate row and column ranks alone
would not suffice.

There are ten consecutive allowed coefficient indices
$b,\ldots,b+9$ with $0\leq b\leq k-9$.
Multiplication by $H$ and restriction to them has determinant
\[
 \det(H_{b+u-v})_{0\leq u,v\leq9}
 =\pm s_{(10^b)}(q^4,q^5,\ldots,q^{k-6})\ne0.
\]
This is the same dual Jacobi--Trudi argument as R-16/R-19
(the cited unconditional identity R-B15). There are $N=k-9$
variables; the numerator exponents are distinct in $[0,N+9]=[0,k]$,
strictly below $m$, so the Weyl alternant is a nonzero Vandermonde.
For $N=0$ the coefficient matrix is the identity.
Consequently the eight normal classes remain independent modulo
the entire base character space after the required coefficient
restriction.

\paragraph{4. Only finitely many modes are lost.}
Choose a nonzero coefficient polynomial in $q$ of the nonzero
polynomial $\mathscr F(q,K,X)$. Its degree is at most $2000$.
For every positive integer $m$,
\[
 \varphi(m)^2/m
 =\prod_{p^e\Vert m}p^{e-2}(p-1)^2\geq\tfrac12:
\]
each odd-prime factor is at least one, and the possible factor
from $p=2$ is at least $1/2$.
For $m\geq M_0$ there are therefore more than $2000$ primitive
$m$th roots. Choose one at which that coefficient polynomial does
not vanish. This uses ordinary root counting and gives
$\mathscr F(q,K,X)\not\equiv0$; it does not incorrectly replace
$\varphi(m)$ by $m$.

For this $q$, at most $300$ of the distinct values $K=q^k$ can
make $\mathscr F(q,K,X)$ the zero polynomial in $X$: choose one
nonzero coefficient polynomial in $K$ and use
\eqref{R20:degrees}. Among $k=9,\ldots,m-9$ there are consequently
at least $m-317$ good bidegrees. For each good $k$, at most $80$
of the distinct values $X=q^t$ are zeros. Thus at least $m-97$
characters in $t=9,\ldots,m-9$ have all eight independent classes.

Consider translates of the eight leading vectors at the $h_m$
block starts $a=5j$, $0\leq j<h_m$. A relation in a good bidegree,
projected onto a good character, gives eight separate equations
\[
 \sum_{j=0}^{h_m-1}c_j(q^{5t})^j=0.
\]
The good characters contain at least $h_m$ distinct values of
$q^{5t}$: every value has at most $\gcd(5,m)\leq5$ preimages,
and $h_m=\lfloor(m-97)/5\rfloor$.
Each equation system contains an ordinary invertible Vandermonde
matrix. Hence the $8h_m$ leading normal classes in each good
bidegree are independent modulo $W_0$.

The original polynomial limits need not be homogeneous for this
bigrading. Nevertheless they are independent modulo $W_0$ as well.
In any relation, take the largest $k$ among vectors with nonzero
coefficients. Their components of bidegree $(m-k,k)$ are precisely
the independent leading classes just proved; every other component
has smaller $(z,w)$ degree. Because $W_0$ is bigraded, its contribution
can be removed in that quotient, forcing all coefficients at this
$k$ to vanish. Descending in $k$ proves independence of every
selected vector. This justifies the leading-component step without
assuming that a projection of the deformed differential image lies
in that image.

Finally, the eight sets of polynomial families all belong, for
nonzero $\varepsilon$, to the differential image of the same
simultaneous pencil. Indeed \eqref{R20:relation} times its outside
permanent is a linear-form combination of global complementary
permanents; division by $\varepsilon^2$ is polynomial by
\eqref{R20:secondorder}. Their limits and $2m^2+2$ base columns
have a nonzero augmented minor at zero. This minor stays nonzero
for a common generic nonzero $\varepsilon$. Thus the dimension
bound \eqref{R20:dimension} follows from the actual differential
rank, with no addition of ranks from different pencils.

\paragraph{5. The uniform padding consequence.}
The general criterion \eqref{R20:allrange} is R-14, with G-10
giving its effective coordinate-degree threshold. Moreover
\[
 h_m\geq(m-101)/5,
\]
so
\begin{align*}
 2m^2+J_m+2-\frac{32m^2}{9}
 &\geq \frac{2m^2}{45}-\frac{3344m}{5}
                  +\frac{256136}{5}+2\;>0
 \qquad(m\geq M_0).
\end{align*}
The last inequality already holds for $m>15048$ after dropping
the positive constant. Thus every integer $m<n<4m/3$ satisfies
\eqref{R20:allrange}. Its strict padding inequality is exactly
the stable range required by G-11, yielding the stated explicit
rectangles. G-12 excludes rectangular gaps for every $n\geq4m/3$.
Finally $J_m=(8/5)m^2+O(m)$, so the full dimension criterion reaches
every fixed ratio below $\sqrt{9/5}=3/\sqrt5$; above $4/3$, the
resulting four-row gaps must be nonrectangular.
\end{proof}

